\section{An optical envelope on shrinking physical segment families}
\label{sec:v87-envelope}
We now use the trace estimate before either the microscopic width or
the collision count is sent to its limit. The source includes every
positive incidence and every value of the optical sensitivity.

For $0<h\le d(a,\delta)$ define
\begin{equation}\label{eq:v87-selected-drift}
 \mathcal V_h^{a,\delta}(t)=
  \sum_j\int_{\mathcal E_t(a,\delta)\cap(\{j\}\times E_j)}
      \int_0^{\min(h,b_j-t)}
         A_j(t+s,\theta)\frac{I_m(t+s,\theta)}{z_j(t+s,\theta)}
                                         \,ds\,d\theta.
\end{equation}
The upper endpoint is interpreted from within the regular segment.
Equivalently this is the reachable-target drift with
$\tau^-(y)>s$, initial gradient $z(\Phi_{-s}y)\ge a$ and total
residence at least $\delta$. It never includes a later re-entry.

\begin{theorem}[A finite-count trace-to-drift estimate]
\label{thm:v87-trace-drift}
For $e_p=1-1/p$ and $0<h\le d(a,\delta)$,
\begin{align}
 \mathcal V_h^{a,\delta}(t)
 &\le \frac{h^{e_p}}a S_t(a,\delta)^{e_p}
                         M_p((t,t+h))^{1/p},\label{eq:v87-holder-drift}\\
 &\le \frac{C_p}a(h/d)^{e_p}M_p(J_t).
                                             \label{eq:v87-linear-drift}
\end{align}
Consequently, for each fixed auxiliary band $B_*$,
\begin{align}
 m^2\mathcal V_h^{a,\delta}(t)
 \le\frac{C_p}a(h/d)^{e_p}\bigl[&C_p(8d+B_*^{-1})\notag\\[-1mm]
 &+C_{B_*,p}(1+8B_*d)m^2
                         \rho_{B_*}^{m(1-p/2)/2}\bigr].
                                           \label{eq:v87-finite-drift}
\end{align}
These inequalities hold at finite count, uniformly in every parameter
in the displayed ranges, in the common essential-roof versions.
\end{theorem}
\begin{proof}
On a selected segment, forward radial monotonicity gives $z(t+s)\ge a$.
Apply H\"older in its time variable, bounding the second factor by
its length to the power $e_p$ times the supremum of $A_j^{e_p}$.
Then apply H\"older to the angular parameter and the disjoint union
of segment charts. This gives \eqref{eq:v87-holder-drift}; the first
factor is precisely the summed trace, not a supremum over words.
The energy in the second factor is a submeasure of the original
$I_m^p$ source on $(t,t+h)$. Insert
Theorem~\ref{thm:v87-summed-trace} and the inclusion
$(t,t+h)\subset J_t$. Since $e_p+1/p=1$, the powers of the larger
energy sum to one, proving \eqref{eq:v87-linear-drift}.
Equation~\eqref{eq:v87-moment} now proves
\eqref{eq:v87-finite-drift}. Nonnegative exhaustion of regular
segments and Tonelli preserve the first-exit selection throughout.
\end{proof}

\begin{theorem}[Uniform microscopic control with vanishing cutoffs]
\label{thm:v87-moving-envelope}
Take $p=33/32$ and set
\[
 a(h)=h^{1/128},\qquad \delta(h)=h^{1/64},\qquad
 c_d=\min\{1/16,(32\Lambda_\circ)^{-1}\}.
\]
There is a fixed $h_*>0$ such that, for every $0<h<h_*$,
$d(a(h),\delta(h))=c_dh^{1/64}\ge h$. For each fixed auxiliary
band there is a finite $K_{B_*}$ satisfying
\begin{equation}\label{eq:v87-uniform-envelope}
 \sup_{m\ge2,R,n,k}m^2
   \|\mathcal V_h^{a(h),\delta(h)}\|_\infty
                         \le K_{B_*}h^{31/1408}.
\end{equation}
No bound on $I_m$, and no internal-incidence corridor, is imposed.
The estimate also bounds, up to a fixed factor, the complete
surviving optical loss and genuine grazing first-exit starting mass
restricted to this initial family. Both restrictions shrink to zero
with $h$ before the collision limit; the estimate is uniform for
all finite collision counts.
\end{theorem}
\begin{proof}
Choose, for example, $h_*<\min\{1,c_d^{64/63}\}$. At fixed $B_*$,
\[
 \sup_{m\ge2}m^2\rho_{B_*}^{m(1-p/2)/2}<\infty.
\]
The bracket in \eqref{eq:v87-finite-drift} is therefore uniformly
bounded for $d\le1/16$. Its outside power of $h$ is exactly
\[
 -\frac1{128}+\frac1{33}\left(1-\frac1{64}\right)
                                  =\frac{31}{1408}>0.
\]
This proves \eqref{eq:v87-uniform-envelope}, with the fixed factor
$c_d^{-1/33}$ included in $K_{B_*}$. One may keep $B_*$ fixed at a
single admissible value for this entire theorem. There is no growing
Fourier band, no infinite sum of collision limsups and no discarded
cutoff-dependent remainder.

For the stated optical consequence, on a surviving segment the
positive decrease of $\omega$ is bounded by its negative variation.
At a genuine grazing first exit its terminal value is zero and the
same negative variation pays its full positive starting weight.
These initial source classes are disjoint. Divide by the initial
$4\pi c_*\chi$, use $W^\varepsilon\le1$ and the two uniform bounds
on $\chi$, and use $|X\omega|\le C\omega I_m/z$. The result is
bounded by $C\mathcal V_h^{a(h),\delta(h)}$. True nongrazing exit
traces remain separate and are not set equal to zero.
\end{proof}

The power $31/1408$ is an explicit admissible choice, not an optimized
exponent. More generally $a(h)=h^\xi$ and $\delta(h)=h^\upsilon$
work whenever $b=\max(2\xi,\upsilon)<1$ and
$e_p(1-b)-\xi>0$. The proved mechanism is physical persistence of
an interior density combined with an actual moment. Unlike the earlier
high-sensitivity envelope, it controls all sensitivities on its stated
geometric family. It does not include source states outside that family.
